% GENERATED FILE. Edit canonical lessons, then run npm run generate:books.
% canonical-source-hash: fnv1a64:66861df132720859
% canonical-lessons: TA-C83-sixth, TA-C83-seventh, TA-C83-eighth, TA-C83-ninth, TA-C83-tenth, TA-C83-eleventh, TA-C83-date, TA-R83-ordinal-recall

\chapter{All the Way Up, and On a Calendar}
\label{ch:ta-ordinals-to-eleventh}
\hlchaptermodality{83}

\begin{chapteropening}
\textbf{By the end of this chapter:} \emph{I can say sixth to eleventh in Tamil, make the ordinal of any number I know, and say a date with the ending a calendar uses.}

Eleven ordinals and two endings retrieved out of order, with three ordinals made for numbers no lesson taught and three dates said with the calendar ending.
\end{chapteropening}

\section[āṟāvadu]{\ta{ஆறாவது} (āṟāvadu) — sixth}

\label{lesson:TA-C83-sixth}



\begin{quote}
\begin{itemize}
\raggedright
  \item \emph{Recall:} say \emph{aindāvadu} — \textbf{R1}
  \item \emph{Recall:} say \emph{iraṇḍāvadu} — \textbf{R2}
  \item \emph{Recall:} give a reason with \emph{ēṉeṉṟāl} — \textbf{R3}
\end{itemize}
\end{quote}

\subsection*{You'll want to know: \ta{ஆறாவது}}
\textbf{\ta{ஆறாவது}} (\emph{āṟāvadu}) — \textquotedblleft{}\textbf{sixth}\textquotedblright{}. \textbf{\ta{ஆறு}} (āṟu) takes \textbf{-\ta{ஆவது}}, and nothing else changes.

Six to ten were a separate lesson from one to five when you learned to count, because the words themselves had to be learned. For ordering they are not a separate anything: the same ending, on the next number.

\subsection*{Your turn}
Say these aloud:

\begin{itemize}
\raggedright
  \item \emph{āṟu}, then \emph{āṟāvadu}
  \item the ordinal before this one, then \emph{āṟāvadu}
  \item \emph{āṟāvadu nāḷ} — the sixth day
\end{itemize}

\begin{tcolorbox}[breakable,colback=teal!4,colframe=teal!35!black,title={Before you move on}]
Did the ending change above five? (\textbf{No}.)

Source: \href{https://www.learntamil.com/part1/ch02/ordinal-numbers-dates/}{learntamil.com, Ordinal Numbers / Dates}.
\end{tcolorbox}
\section[ēḻāvadu]{\ta{ஏழாவது} (ēḻāvadu) — seventh}

\label{lesson:TA-C83-seventh}



\begin{quote}
\begin{itemize}
\raggedright
  \item \emph{Recall:} say \emph{āṟāvadu} — \textbf{R1}
  \item \emph{Recall:} say \emph{mutalāvadu} — \textbf{R2}
  \item \emph{Recall:} read \textbf{\ta{ஏன்}} — \textbf{R3}
  \item \emph{Recall:} say \emph{saṭṭai} — \textbf{R4}
\end{itemize}
\end{quote}

\subsection*{You'll want to know: \ta{ஏழாவது}}
\textbf{\ta{ஏழாவது}} (\emph{ēḻāvadu}) — \textquotedblleft{}\textbf{seventh}\textquotedblright{}. \textbf{\ta{ஏழு}} (ēḻu) takes \textbf{-\ta{ஆவது}}, and nothing else changes.

\textbf{\ta{ஏழு}} carries the curled \textbf{ḻ}, the sound in the name \emph{Tamiḻ} itself and the one this language shares with Malayalam while Kannada and Telugu hardened it away. Adding the ending does not soften it: \textbf{\ta{ஏழாவது}}, tongue still curled.

\subsection*{Your turn}
Say these aloud:

\begin{itemize}
\raggedright
  \item \emph{ēḻu}, then \emph{ēḻāvadu}
  \item the ordinal before this one, then \emph{ēḻāvadu}
  \item \emph{ēḻāvadu nāḷ} — the seventh day
\end{itemize}

\begin{tcolorbox}[breakable,colback=teal!4,colframe=teal!35!black,title={Before you move on}]
What happens to the \ta{ழ} when the ending goes on? (\textbf{Nothing}.)

Source: \href{https://www.learntamil.com/part1/ch02/ordinal-numbers-dates/}{learntamil.com, Ordinal Numbers / Dates}.
\end{tcolorbox}
\section[eṭṭāvadu]{\ta{எட்டாவது} (eṭṭāvadu) — eighth}

\label{lesson:TA-C83-eighth}



\begin{quote}
\begin{itemize}
\raggedright
  \item \emph{Recall:} say \emph{ēḻāvadu} — \textbf{R1}
  \item \emph{Recall:} say \emph{mūṉṟāvadu} — \textbf{R2}
  \item \emph{Recall:} say \emph{adaṉāl} — \textbf{R3}
  \item \emph{Recall:} say \emph{vēṭṭi} — \textbf{R4}
\end{itemize}
\end{quote}

\subsection*{You'll want to know: \ta{எட்டாவது}}
\textbf{\ta{எட்டாவது}} (\emph{eṭṭāvadu}) — \textquotedblleft{}\textbf{eighth}\textquotedblright{}. \textbf{\ta{எட்டு}} (eṭṭu) takes \textbf{-\ta{ஆவது}}, and nothing else changes.

\textbf{\ta{எட்டு}} has the doubled retroflex \textbf{ṭṭ}. The ending lands after all of it, so the doubling is untouched: \textbf{\ta{எட்டாவது}}.

\subsection*{Your turn}
Say these aloud:

\begin{itemize}
\raggedright
  \item \emph{eṭṭu}, then \emph{eṭṭāvadu}
  \item the ordinal before this one, then \emph{eṭṭāvadu}
  \item \emph{eṭṭāvadu nāḷ} — the eighth day
\end{itemize}

\begin{tcolorbox}[breakable,colback=teal!4,colframe=teal!35!black,title={Before you move on}]
Did the doubled \ta{ட்ட} change? (\textbf{No}.)

Source: \href{https://www.learntamil.com/part1/ch02/ordinal-numbers-dates/}{learntamil.com, Ordinal Numbers / Dates}.
\end{tcolorbox}
\section[oṉpadāvadu]{\ta{ஒன்பதாவது} (oṉpadāvadu) — ninth}

\label{lesson:TA-C83-ninth}



\begin{quote}
\begin{itemize}
\raggedright
  \item \emph{Recall:} say \emph{eṭṭāvadu} — \textbf{R1}
  \item \emph{Recall:} say \emph{nāṉkāvadu} — \textbf{R2}
  \item \emph{Recall:} say why you are doing something with \emph{-kkāga} — \textbf{R3}
  \item \emph{Recall:} read \textbf{\ta{சட்டை}} — \textbf{R4}
\end{itemize}
\end{quote}

\subsection*{You'll want to know: \ta{ஒன்பதாவது}}
\textbf{\ta{ஒன்பதாவது}} (\emph{oṉpadāvadu}) — \textquotedblleft{}\textbf{ninth}\textquotedblright{}. \textbf{\ta{ஒன்பது}} (oṉpadu) takes \textbf{-\ta{ஆவது}}, and nothing else changes.

\textbf{\ta{ஒன்பது}} was already a built word when you learned it: an old \textbf{\ta{ஒன்}}, one, in front of \textbf{\ta{பது}}, ten — one short of ten. The ordinal ending goes on top of that, and none of the machinery underneath reacts to it.

\subsection*{Your turn}
Say these aloud:

\begin{itemize}
\raggedright
  \item \emph{oṉpadu}, then \emph{oṉpadāvadu}
  \item the ordinal before this one, then \emph{oṉpadāvadu}
  \item \emph{oṉpadāvadu nāḷ} — the ninth day
\end{itemize}

\begin{tcolorbox}[breakable,colback=teal!4,colframe=teal!35!black,title={Before you move on}]
What two numbers are inside \ta{ஒன்பது}? (\textbf{One} and \textbf{ten} — nine is one short of ten.)

Source: \href{https://www.learntamil.com/part1/ch02/ordinal-numbers-dates/}{learntamil.com, Ordinal Numbers / Dates}.
\end{tcolorbox}
\section[pattāvadu]{\ta{பத்தாவது} (pattāvadu) — tenth}

\label{lesson:TA-C83-tenth}



\begin{quote}
\begin{itemize}
\raggedright
  \item \emph{Recall:} say \emph{oṉpadāvadu} — \textbf{R1}
  \item \emph{Recall:} say \emph{aindāvadu} — \textbf{R2}
  \item \emph{Recall:} ask \emph{eppōdu?} — \textbf{R3}
  \item \emph{Recall:} say \emph{puḍavai} — \textbf{R4}
\end{itemize}
\end{quote}

\subsection*{You'll want to know: \ta{பத்தாவது}}
\textbf{\ta{பத்தாவது}} (\emph{pattāvadu}) — \textquotedblleft{}\textbf{tenth}\textquotedblright{}. \textbf{\ta{பத்து}} (pattu) takes \textbf{-\ta{ஆவது}}, and nothing else changes.

\textbf{\ta{பத்தாவது} \ta{நாள்}} (\emph{pattāvadu nāḷ}), the tenth day, and the ten are done.

\subsection*{Your turn}
Say these aloud:

\begin{itemize}
\raggedright
  \item \emph{pattu}, then \emph{pattāvadu}
  \item the ordinal before this one, then \emph{pattāvadu}
  \item \emph{pattāvadu nāḷ} — the tenth day
\end{itemize}

\begin{tcolorbox}[breakable,colback=teal!4,colframe=teal!35!black,title={Before you move on}]
Count the ordinals from one to ten. How many of them were words you had to learn? (\textbf{None} — ten of them were built by one ending on numbers you already count with, and \ta{முதல்} you already had.)

Source: \href{https://www.learntamil.com/part1/ch02/ordinal-numbers-dates/}{learntamil.com, Ordinal Numbers / Dates}.
\end{tcolorbox}
\section[patiṉoṉṟāvadu]{\ta{பதினொன்றாவது} (patiṉoṉṟāvadu) — eleventh, and the edge that is not there}

\label{lesson:TA-C83-eleventh}



\begin{quote}
\begin{itemize}
\raggedright
  \item \emph{Recall:} say \emph{pattāvadu} — \textbf{R1}
  \item \emph{Recall:} say \emph{āṟāvadu} — \textbf{R2}
  \item \emph{Recall:} say \emph{appōdu} — \textbf{R3}
  \item \emph{Recall:} say \emph{seruppu} — \textbf{R4}
\end{itemize}
\end{quote}

\subsection*{You'll want to know: \ta{பதினொன்றாவது}}
\textbf{\ta{பதினொன்றாவது}} (\emph{patiṉoṉṟāvadu}) — \textquotedblleft{}\textbf{eleventh}\textquotedblright{}.

\textbf{\ta{பதினொன்று}} was itself a built word: \emph{patiṉ-}, the joining shape of ten, with \emph{oṉṟu} welded on. Now the ordinal ending goes on top of \emph{that}, and the result is a word of three parts, none of which had to be learned as a whole.

\begin{grammarlens}[title={Grammar Lens: the ending is a fact about numbers, not about ten}]
Most books stop an ordinal chapter at \emph{tenth}, and a learner reasonably concludes that those ten words were the set. In Tamil there is no set. \textbf{-\ta{ஆவது}} attaches to a \textbf{number}, so every number you learn from here on arrives with its ordinal already made: \textbf{\ta{இருபது}} gives \textbf{\ta{இருபதாவது}} (\emph{irupadāvadu}), twentieth, without a lesson.

That is why this book taught the ending in the second lesson of the previous chapter rather than after ten examples. The examples were never the content.
\end{grammarlens}

\subsection*{Your turn}
Say these aloud:

\begin{itemize}
\raggedright
  \item \emph{patiṉoṉṟu}, then \emph{patiṉoṉṟāvadu}
  \item \emph{pattāvadu, patiṉoṉṟāvadu}
  \item take \emph{irupadu} and make its ordinal yourself
\end{itemize}

\begin{tcolorbox}[breakable,colback=teal!4,colframe=teal!35!black,title={Before you move on}]
Do you need a lesson for \emph{twentieth}? (\textbf{No} — the ending goes on any number.)

Source: \href{https://www.learntamil.com/part1/ch02/ordinal-numbers-dates/}{learntamil.com, Ordinal Numbers / Dates}.
\end{tcolorbox}
\section[tēdi]{\ta{தேதி} (tēdi) — the date, and the ending a calendar uses}

\label{lesson:TA-C83-date}



\begin{quote}
\begin{itemize}
\raggedright
  \item \emph{Recall:} say \emph{patiṉoṉṟāvadu}, and say why \emph{irupadāvadu} needs no lesson — \textbf{R1}
  \item \emph{Recall:} say \emph{ēḻāvadu} — \textbf{R2}
  \item \emph{Recall:} read \textbf{\ta{எப்போது}} — \textbf{R3}
  \item \emph{Recall:} say \emph{toppi} — \textbf{R4}
\end{itemize}
\end{quote}

\subsection*{You'll want to know: \ta{தேதி}}
\textbf{\ta{தேதி}} (\emph{tēdi}) — \textquotedblleft{}\textbf{a date}\textquotedblright{}, the day-number on a calendar.

\begin{grammarlens}[title={Grammar Lens: a second ending, for one job}]
Tamil has a \textbf{second} ordinal ending, and it is not a free alternative to the first: \textbf{-\ta{ஆம்}} (\emph{-ām}) is what a \textbf{date} takes.

\noindent
\begin{tabularx}{\linewidth}{@{}>{\raggedright\arraybackslash}X>{\raggedright\arraybackslash}X@{}}
\toprule
\textbf{said of a place in a line} & \textbf{said of a date} \\
\midrule
\textbf{\ta{இரண்டாவது}} \emph{iraṇḍāvadu} & \textbf{\ta{இரண்டாம்} \ta{தேதி}} \emph{iraṇḍām tēdi} \\
\textbf{\ta{மூன்றாவது}} \emph{mūṉṟāvadu} & \textbf{\ta{மூன்றாம்} \ta{தேதி}} \emph{mūṉṟām tēdi} \\
\bottomrule
\end{tabularx}

Same number, same job in the world, two endings — and the calendar keeps the older, shorter one.

You already have the months. \textbf{\ta{சித்திரை} \ta{இரண்டாம்} \ta{தேதி}} (\emph{Chithirai iraṇḍām tēdi}) is the second of Chithirai, and that is a date you can now write down.
\end{grammarlens}

\subsection*{Your turn}
Say these aloud:

\begin{itemize}
\raggedright
  \item \emph{iraṇḍāvadu}, then \emph{iraṇḍām tēdi}, and say which one is the date
  \item \emph{Chithirai iraṇḍām tēdi}
  \item today's day-number with the ending a calendar uses
\end{itemize}

\begin{tcolorbox}[breakable,colback=teal!4,colframe=teal!35!black,title={Before you move on}]
Which ending does a date take, \textbf{-\ta{ஆவது}} or \textbf{-\ta{ஆம்}}? (\textbf{-\ta{ஆம்}}.) Say \textquotedblleft{}the third of the month\textquotedblright{}. (\textbf{\ta{மூன்றாம்} \ta{தேதி}}.)

Sources: \href{https://www.learntamil.com/part1/ch02/ordinal-numbers-dates/}{learntamil.com, Ordinal Numbers / Dates}; \href{https://www.omniglot.com/language/numbers/tamil.htm}{Omniglot, Tamil numbers}.
\end{tcolorbox}
\section[Practice]{The ordinals, at a distance}

\label{lesson:TA-R83-ordinal-recall}



\begin{quote}
\begin{itemize}
\raggedright
  \item \emph{Recall:} say \emph{iraṇḍām tēdi}, and say which ending a date takes — \textbf{R1}
  \item \emph{Recall:} say \emph{eṭṭāvadu} — \textbf{R2}
  \item \emph{Recall:} build a when-clause with \emph{-um pōdu} — \textbf{R3}
  \item \emph{Recall:} say \emph{kaḍai} — \textbf{R4}
\end{itemize}
\end{quote}

\subsection*{Your turn}
Eleven ordinals, two endings and one word that was never new. Run it four ways.

\begin{enumerate}
\raggedright
  \item Say first to eleventh straight through, once slowly and once fast.
  \item Hear six of them out of order; say each back and say the one above it.
  \item Take three numbers you were \textbf{not} taught an ordinal for and make theirs.
  \item Say three dates with the month names you already have, using the ending a calendar takes rather than the ordinary one.
\end{enumerate}

Then answer aloud: which of the eleven is not built by the ending on its own number, what is it built on instead, and where in this book did you first say it?

\begin{tcolorbox}[breakable,colback=teal!4,colframe=teal!35!black,title={Before you move on}]
Which word for \textbf{first} was already inside a word you knew? (\textbf{\ta{முதல்}}, inside \textbf{\ta{முதலில்}}.) Which ending does a date take? (\textbf{-\ta{ஆம்}}.)

Sources: \href{https://www.learntamil.com/part1/ch02/ordinal-numbers-dates/}{learntamil.com, Ordinal Numbers / Dates}; \href{https://www.omniglot.com/language/numbers/tamil.htm}{Omniglot, Tamil numbers}.
\end{tcolorbox}
